\documentclass[slideColor]{prosper}
\usepackage{amsmath,amsfonts,amsthm}
%\usepackage{epsfig}
%\include{Qcircuit}


\newcommand{\bbC}{{\mathbb{C}}}
\newcommand{\C}{{\mathbb{C}}}
\newcommand{\F}{{\mathbb{F}}}
\newcommand{\N}{{\mathbb{N}}}
\newcommand{\Q}{{\mathbb{Q}}}
\renewcommand{\R}{{\mathbb{R}}}
\newcommand{\Z}{{\mathbb{Z}}}

\newcommand{\cD}{{\mathcal{D}}}
\newcommand{\cN}{{\mathcal{N}}}
\newcommand{\cP}{{\mathcal{P}}}
\newcommand{\cQ}{{\mathcal{Q}}}
\newcommand{\cS}{{\mathcal{S}}}
\newcommand{\cU}{{\mathcal{U}}}
\newcommand{\cV}{{\mathcal{V}}}

\newcommand{\tr}{\mathop{\mathrm{tr}}}
\newcommand{\rank}{\mathop{\mathrm{rank}}}
\newcommand{\poly}{\mathop{\mathrm{poly}}}
\newcommand{\E}{\mathop{\mbox{$\mathbf{E}$}}}
\newcommand{\schur}{\mathop{\mathrm{Schur}}}
\newcommand{\schursmall}[2]{{\mathrm{S}}_{#1,#2}}
\newcommand{\planch}{\mathop{\mathrm{Planch}}}
\newcommand{\planchsmall}[1]{{\mathrm{P}}_{#1}}
\renewcommand{\d}{{\mathrm{d}}}

\newcommand{\FF}{{\mathbb F}}
\newcommand{\ZZ}{{\mathbb Z}}
\newcommand{\QQ}{{\mathbb Q}}
\newcommand{\NN}{{\mathbb N}}
\renewcommand{\AA}{{\mathbb A}}


\newcommand{\HSP}{\textsc{hsp}\xspace}

\renewcommand{\>}{\rangle}
\newcommand{\<}{\langle}
\newcommand{\ket}[1]{|#1\rangle}
\newcommand{\bra}[1]{\langle #1|}
\newcommand{\scalar}[2]{\langle #1|#2\rangle}
\newcommand{\proj}[1]{\left|#1\right\>\!\left\<#1\right|}
\newcommand{\norm}[1]{\|#1\|}
\newcommand{\ot}{\otimes}
\newcommand{\eps}{\epsilon}
\newcommand{\ra}{\rightarrow}

\renewcommand{\setminus}{\mathbin{-}}
\newcommand{\partitionof}{\mathbin{\vdash}}


%%%%%%%%%%%%%%%%%%%%%%%%%%%%%%%%%%%%%%%%%%%%%%%%%%%%%%%%%%%%

\title{Efficient Quantum Algorithm for Identifying Hidden Polynomial Function Graphs}
\author{Pawel M. Wocjan\\[1ex]School of Electrical Engineering and Computer Science\\ University of Central Florida\\Orlando}
\email{wocjan@cs.ucf.edu}

\begin{document}

\maketitle

%%%%%%%%%%%%%%%%%%%%%%%%%%%%%%%%%%%%%%%%%%%%%%%%%%%%%%%%%%%%%%%%%%%%%%%%

\begin{slide}{Joint work}
\begin{itemize}
\item Thomas Decker (Department of Computer Science \& Engineering, University of Washington)
\item Jan Draisma (Department of Mathematics \& Computer Science, Technical University of Eindhoven)
\end{itemize}
\end{slide}

%%%%%%%%%%%%%%%%%%%%%%%%%%%%%%%%%%%%%%%%%%%%%%%%%%%%%%%%%%%%%%%%%%%%%%%%

\begin{slide}{Motivation and Context}

\begin{itemize}
\item Which problems can be solved exponentially faster on quantum computer than on a classical computer?
\item Most such problems are black-box problems
\begin{itemize}
\item Hidden Subgroup Problems (abelian/nonabelian)
\item Hidden Shift Problems
\end{itemize}
\item Motivation: interesting problems can be reduced to these black-box problems
\begin{itemize}
\item Integer Factorization, Discrete Logarithm
\item Pell's Equation, Principle Ideal Problem, Problem of Computing the Unit and Class Group of a Number Field
\item Shortest Lattice Problem
\item Graph Isomorphism
\end{itemize}
\end{itemize}
\end{slide}

%%%%%%%%%%%%%%%%%%%%%%%%%%%%%%%%%%%%%%%%%%%%%%%%%%%%%%%%%%%%%%%%%%%%%%%%

\begin{slide}{Motivation and Context}
\begin{itemize}
\item Most quantum algorithms rely on the efficient solution of the abelian HSP
\item The nonabelian HSP seems to be hard for groups of interest such as the symmetric and dihedral groups
\item Are there other ways of extending the success of the abelian HSP besides considering the nonabelian case?
\end{itemize}
\end{slide}

\begin{slide}{Non-linear Structures}
Design quantum algorithms for certain black-box problems involving \textcolor{green}{non-linear structures}

\bigskip
A.~Childs, L.~Schulman, and U.~Vazirani, {\em Quantum algorithms for hidden nonlinear structures}, QIP 2007
\begin{itemize}
\item Query complexity of Hidden Polynomial Problem
\item Efficient quantum algorithms for Hidden Sphere Problems
\begin{itemize}
\item Hidden Radius Problem
\item Hidden Subspace of Centers Problem
\end{itemize}
\end{itemize}
\end{slide}

%%%%%%%%%%%%%%%%%%%%%%%%%%%%%%%%%%%%%%%%%%%%%%%%%%%%%%%%%%%%%%%%%%%%%%%%

\begin{slide}{Outline}
\begin{itemize}
\item Definition of the Hidden Polynomial Function Graph Problem
\item Reduction to quantum state identification
\item Polynomial quantum query complexity
\item Construction of a von Neuman measurement for identifying the quantum states
\begin{itemize}
\item Reduction of its implementation to a classical algebro-geometric problem
\item Efficient solution of this problem
\end{itemize}
\end{itemize}
\end{slide}

%%%%%%%%%%%%%%%%%%%%%%%%%%%%%%%%%%%%%%%%%%%%%%%%%%%%%%%%%%%%%%%%%%%%%%%%

\begin{slide}{Definition}
\begin{itemize}
\item let $Q(X)\in\F[X]$ be an arbitrary polynomial of total degree at most $n$ whose constant term is equal to $0$
\item let $B : \F^2 \rightarrow \F$ be a black-box function hiding the polynomial Q in the following sense:
\[
B(x,y)=\pi(y-Q(x))\,,
\]
where $\pi$ is an unknown bijection permuting the elements of $\F$ arbitrarily
\item the Hidden Polynomial Function Graph Problem is to identify $Q$
\item an algorithm for polynomials with degree less or equal to $n$ (where $n$ is a constant) is efficient if its running time is polylogarithmic in $p$, i.e., the size of $\F$
\end{itemize}
\end{slide}

%%%%%%%%%%%%%%%%%%%%%%%%%%%%%%%%%%%%%%%%%%%%%%%%%%%%%%%%%%%%%%%%%%%%%%%%

\begin{slide}{Hidden Polynomial Function Graph}
can be considered as a natural generalization of abelian HSP over groups of form $\F\times\F$
\begin{itemize}
\item the black-box function $B$ is constant on the subsets
\[
H_{Q,z}:=\{(x,Q(x)+z) : x \in \F\}
\]
of $\F^2$ and distinct for different values of $z$
\item assume that $Q$ is linear, i.e., $Q(X)=q X$ with unknown $q$ $\Rightarrow$
\[
H_{Q}:=\{(x,Q(x)) : x\in\F\}
\] 
is a subgroup as $(x,Q(x))+(x',Q(x'))=(x+x',Q(x,x'))$
\item the sets $H_{Q,z}$ correspond to the cosets of the hidden subgroup $H_Q$
\end{itemize}
\end{slide}

%%%%%%%%%%%%%%%%%%%%%%%%%%%%%%%%%%%%%%%%%%%%%%%%%%%%%%%%%%%%%%%%%%%%%%%%

\begin{slide}{Classical Query Complexity}
\begin{itemize}
\item we need to find at least $n$ different points 
\[
(x_1,y_1), \ldots (x_n,y_n) \mbox{ with } B(x_1,y_1)=\ldots B(x_n,y_n)
\]
in order to determine the hidden polynomial $Q$ of degree $n$
\item the probability of obtaining such an $n$-fold collision is $1/p^{n-1}$
\item $\Rightarrow$ classical query complexity is polynomial in $p$
\end{itemize}
\end{slide}

%%%%%%%%%%%%%%%%%%%%%%%%%%%%%%%%%%%%%%%%%%%%%%%%%%%%%%%%%%%%%%%%%%%%%%%%

\begin{slide}{Reduction to Quantum State Identification}
\begin{itemize}
\item evaluate the black-box function on an equally weighted superposition of $(x,y)\in\F^2$
\[
\frac{1}{p} \sum_{x,y\in\F} |x\> \otimes |y\> \otimes |B(x,y)\>
\]
\item measure third register; assume we have obtained $\pi(z)$, then the state on the first and second register is
\[
|\phi_{Q,z}\> := \frac{1}{\sqrt{p}}\sum_{x\in\F} |x\> \otimes |Q(x)+z\>\,,
\]
where $z\in\F$ is uniformly at random; the density matrix $\rho_Q := \frac{1}{p}\sum_{z\in\F} |\phi_{Q,z}\>\<\phi_{Q,z}|$ is called a polynomial function state
\end{itemize}
\end{slide}

%%%%%%%%%%%%%%%%%%%%%%%%%%%%%%%%%%%%%%%%%%%%%%%%%%%%%%%%%%%%%%%%%%%%%%%%

\begin{slide}{Quantum Query Complexity}
\begin{itemize}
\item let $\rho$ and $\sigma$ be two quantum state with spectral decompositions $\rho=\sum_i \lambda_i |\psi_i\>\<\psi_i|$ and $\rho=\sum_i \mu_j |\phi_j\>\<\phi_j|$
\item assume that $\max_{i,j}|\<\psi_i|\phi_j\>|\le\alpha$ for some $\alpha$ $\Rightarrow$
\[
F(\rho,\sigma)\le\alpha \cdot \min \left\{
\sum_i \sqrt{\lambda_i}, \sqrt{\mu_j}
\right\}\,,
\]
where $F(\sigma,\rho):=\|\sqrt{\rho}\sqrt{\sigma}\|_1$ is the fidelity of $\rho$ and $\sigma$
\item we have $F(\rho_Q,\rho_{\tilde{Q}})\le n/\sqrt{p}$ for any pair of different polynomial function states as
\[
|\<\phi_{Q,z}|\phi_{\tilde{Q},\tilde{z}}\> = \frac{1}{p} |\{x \in \F : Q(x) + z = \tilde{Q}(x) + \tilde{z}\}| \le
\frac{n}{p}
\]
\end{itemize}
\end{slide}

%%%%%%%%%%%%%%%%%%%%%%%%%%%%%%%%%%%%%%%%%%%%%%%%%%%%%%%%%%%%%%%%%%%%%%%%

\begin{slide}{Quantum Query Complexity}
\begin{itemize}
\item there is a POVM $\{E_Q\}$ acting on $k$ copies of a polynomial function state such that
\[
P_{\rm success} := \min_Q \mathrm{Tr} (\rho_Q^{\otimes k} E_Q) \ge 1-\epsilon
\]
provided that $k\ge 2(\log N - \log\epsilon)/(-\log F)$, where 
\begin{itemize}
\item $N:=p^n$ is the number of different $\rho_Q$ 
\item $F:=\max_{Q\neq\tilde{Q}} F(\rho_Q,\rho_{\tilde{Q}})$
\end{itemize}
(this follows from results by {\sc Harrow and Winter})
\item this bound and the upper bound on the maximal fidelity $F\le n/\sqrt{p}$ imply $P_{\rm success}$ is a least $1/2$ for $k=4n$
\item $n$ copies are required to have $P_{\rm success}$ at least $1/2$ (follows from a lower bound)
\end{itemize}
\end{slide}

%%%%%%%%%%%%%%%%%%%%%%%%%%%%%%%%%%%%%%%%%%%%%%%%%%%%%%%%%%%%%%%%%%%%%%%%

\begin{slide}{Structure of $\rho_Q$}
\begin{itemize}
\item the state $\rho_Q$ can be written as
\begin{eqnarray*}
\rho_Q & = & 
\frac{1}{p^2} \sum_{b,c\in\F} |b\>\<c| \otimes \sum_{z\in\F} |Q(b)+z\>\<Q(c)+z| \\
& = & 
\frac{1}{p^2} \sum_{b,c\in\F} \sum|b\>\<c| \otimes S^{Q(b)-Q(c)}\,,
\end{eqnarray*}
where $S$ is the cyclic shift operator acting by $S|u\>=|u+1 \mod p\>$
\item apply the Fourier transform $F$ on the second register to diagonalize all powers of $S$
\[
\hat{\rho}_Q := 
({\bf 1} \otimes F) \rho ({\bf 1} \otimes F)^\dagger
=
\frac{1}{p^2} \sum_{b,c,x\in\F} \omega_p^{[Q(b)-Q(c)]\cdot x} |b\>\<c| \otimes |x\>\<x|
\]
\end{itemize}
\end{slide}

%%%%%%%%%%%%%%%%%%%%%%%%%%%%%%%%%%%%%%%%%%%%%%%%%%%%%%%%%%%%%%%%%%%%%%%%

\begin{slide}{Structure of $\hat{\rho}_Q^{\otimes k}$}
\vspace{-0.5cm}
\begin{eqnarray*}
\hat{\rho}_Q^{\otimes k} & = & 
\frac{1}{p^{2k}}
\sum_{b,c,x\in\F^k} \omega_p^{\sum_{j=1}^k [Q(b_j)-Q(c_j)] x_j}
|b\>\<c| \otimes |x\>\<x| \\
& = &
\frac{1}{p^{2k}} \sum_{b,c,x\in\F^k}
\omega_p^{\<q|\Phi_n(b) - \Phi_n(c)|x\>} |b\>\<c| \otimes
|x\>\<x|
\end{eqnarray*}
\begin{itemize}
\item $\<q|:=(q_1,q_2,\ldots,q_n) \in \F^{1 \times n}$ where $Q(X)=\sum_{i=1}^n q_i X^i$
\item 
\[
\Phi_n(b) := \sum_{i=1}^n \sum_{j=1}^k b_j^i |i\>\<j| = \left(
\begin{array}{cccc}
b_1    & b_2    & \cdots & b_k    \\
b_1^2  & b_2^2  & \cdots & b_k^2  \\
\vdots & \vdots &        & \vdots \\
b_1^n  & b_2^n  & \cdots & b_k^n  
\end{array}
\right)\,,
\]
\item \vspace{-0.2cm} $|x\>:=(x_1,\ldots,x_k)^T \in \F^k$
\end{itemize}
\end{slide}

%%%%%%%%%%%%%%%%%%%%%%%%%%%%%%%%%%%%%%%%%%%%%%%%%%%%%%%%%%%%%%%%%%%%%%%%

\begin{slide}{Algebro-Geometric Problem}
\vspace{-0.2cm}
\begin{itemize}
\item determine all $b\in\F^k$ for given $x\in\F^k$ and $w\in\F^n$ s.t.
\[
\left(
\begin{array}{cccc}
b_1 & b_2 & \cdots & b_k \\ b_1^2 & b_2^2 & \cdots & b_k^2 \\ \vdots &
\vdots & & \vdots \\ b_1^n & b_2^n & \cdots & b_k^n
\end{array}
\right)\cdot \left(
\begin{array}{c}
x_1 \\ x_2 \\ \vdots \\ x_k
\end{array}
\right) = \left(
\begin{array}{c}
w_1 \\ w_2 \\ \vdots \\ w_n
\end{array}
\right)
\]
\item denote the set of solutions and its cardinality by
\[
S_w^x := \{ b\in \F^k \,:\, \Phi_n(b)|x\> = |w\>\} 
\quad\; {\rm and} \quad\;
\eta_w^x := |S_w^x|
\]
\item define the quantum states $|S_w^x\>$ for $\eta^x_w>0$ 
\[
|S_w^x\> := \frac{1}{\sqrt{\eta_w^x}} \sum_{b\in S_w^x} |b\>
\]
\end{itemize}
\end{slide}

%%%%%%%%%%%%%%%%%%%%%%%%%%%%%%%%%%%%%%%%%%%%%%%%%%%%%%%%%%%%%%%%%%%%%%%%

\begin{slide}{Structure of $\hat{\rho}_Q^{\otimes k}$}
\begin{eqnarray*}
\hat{\rho}_Q^{\otimes k} & = & 
\frac{1}{p^{2k}} \sum_{\textcolor{red}{b},\textcolor{blue}{c},x\in\F^k}
\omega_p^{\<q| \textcolor{red}{\Phi_n(b)|x\>} - \<q|\textcolor{blue}{\Phi_n(c)|x\>}} 
|\textcolor{red}{b}\>\<\textcolor{blue}{c}| \otimes |x\>\<x| \\
& = &
\frac{1}{p^{2k}}\sum_{x\in\F^k}
\sum_{\textcolor{red}{w},\textcolor{blue}{v}\in\F^n} \omega_p^{\<q|\textcolor{red}{w\>} - \<q|\textcolor{blue}{v\>}} 
\sqrt{\textcolor{red}{\eta_w^x} \textcolor{blue}{\eta_v^x}} 
|\textcolor{red}{S_w^x}\>\<\textcolor{blue}{S_v^x}| \otimes |x\>\<x|
\end{eqnarray*}
\end{slide}

%%%%%%%%%%%%%%%%%%%%%%%%%%%%%%%%%%%%%%%%%%%%%%%%%%%%%%%%%%%%%%%%%%%%%%%%

\begin{slide}{Measurement of $x$}
\begin{itemize}
\item the block structure implies that we can measure $x\in\F^k$ without loss of generality
\item the probability of obtaining a particular $x$ is $1/p^k$ (uniformly at random)
\item the reduced state is
\[
\hat{\rho}_Q^x := 
\frac{1}{p^{k}}
\sum_{w,v\in\F^n} \omega_p^{\<q|w\> - \<q|v\>} \sqrt{\eta_w^x
\eta_v^x} |S_w^x\>\<S_v^x| 
\]
\end{itemize}
\end{slide}

%%%%%%%%%%%%%%%%%%%%%%%%%%%%%%%%%%%%%%%%%%%%%%%%%%%%%%%%%%%%%%%%%%%%%%%%

\begin{slide}{Ideal Quantum Sampling}
\begin{itemize}
\item assume we could implement a unitary $U_x$ with
\[
U_x |S_w^x\> = |w\>
\]
\item applying $U_x$ to $\hat{\rho}_Q^x$ gives
\[
U_x \hat{\rho}_Q^x U_x^\dagger = \frac{1}{p^k} \sum_{w,v\in\F^n}
\omega_p^{\<q|w\> - \<q|v\>} \sqrt{\eta_w^x \eta_v^x} |w\>\<v|
\]
\item
\end{itemize}
\end{slide}

%%%%%%%%%%%%%%%%%%%%%%%%%%%%%%%%%%%%%%%%%%%%%%%%%%%%%%%%%%%%%%%%%%%%%%%%

\begin{slide}{Measurement in Fourier basis}
\begin{itemize}
\item measure in the Fourier basis, i.e., carry out the orthogonal measurement with respect to the states
\[
|\psi_Q\> := \frac{1}{\sqrt{p^n}} \sum_{w \in \F^n}
 \omega_p^{\<q|w\>} |w\>
\]
\item the probability for correct detection of $\hat{\rho}_Q^x$ is
\[
\< \psi_Q | U_x {\hat \rho}_Q^x U_x^\dagger | \psi_Q \> =
\frac{1}{p^{k+n}} \left( \sum_{w \in \F^n} \sqrt{\eta_w^x}
\right)^2\
\]
\end{itemize}
\end{slide}

%%%%%%%%%%%%%%%%%%%%%%%%%%%%%%%%%%%%%%%%%%%%%%%%%%%%%%%%%%%%%%%%%%%%%%%%

\begin{slide}{Summery of the Steps}
\begin{itemize}
\item apply Fourier on second register
\item measure $x$ in second register
\item apply (ideal) quantum sampling $U_x$
\item measure in Fourier basis
\item the outcome determines the coefficients of the hidden polynomial
\item the overall success probability is
\[
\frac{1}{p^{2k+n}} \sum_{x\in\F^k} \left( \sum_{w \in \F_p^n} \sqrt{\eta_w^x}
\right)^2
\]
\end{itemize}
\end{slide}

%%%%%%%%%%%%%%%%%%%%%%%%%%%%%%%%%%%%%%%%%%%%%%%%%%%%%%%%%%%%%%%%%%%%%%%%

\begin{slide}{Approximate Quantum Sampling}
\vspace{-0.2cm}
\begin{itemize}
\item determine all $b\in\F^k$ for given $x\in\F^k$ and $w\in\F^n$ 
\[
\left(
\begin{array}{cccc}
b_1 & b_2 & \cdots & b_k \\ b_1^2 & b_2^2 & \cdots & b_k^2 \\ \vdots &
\vdots & & \vdots \\ b_1^n & b_2^n & \cdots & b_k^n
\end{array}
\right)\cdot \left(
\begin{array}{c}
x_1 \\ x_2 \\ \vdots \\ x_k
\end{array}
\right) = \left(
\begin{array}{c}
w_1 \\ w_2 \\ \vdots \\ w_n
\end{array}
\right)
\]
\item recall $S_w^x := \{ b\in \F^k \,:\, \Phi_n(b)|x\> = |w\>\}$ and $\eta_w^x := |S_w^x|$
\item it suffices to implement an approximate version $V_x$ of quantum sampling satisfying
\[
V_x|S_w^x\>=|w\>
\]
for a vast majority of pairs $(x,w)$.
\end{itemize}
\end{slide}

%%%%%%%%%%%%%%%%%%%%%%%%%%%%%%%%%%%%%%%%%%%%%%%%%%%%%%%%%%%%%%%%%%%%%%%%

\begin{slide}{Properties of Generic Morphisms}
Proposition
\begin{itemize}
\item
consider a morphism $f:\AA^k \times \AA^n \rightarrow \AA^n$ defined
over $\ZZ$
\begin{itemize}
\item $f$ is given by an $n$-tuple $f=(f_1,\ldots,f_n)$
of polynomials in $\ZZ[X,B]$
\item $X=(X_1,\ldots,X_k)$ and
$B=(B_1,\ldots,B_n)$ are the coordinates on $\AA^k$ and on the first
copy of $\AA^n$ 
\end{itemize}
\item 
suppose that the Jacobian determinant
$\det (\partial f_i/\partial B_j)_{ij}$ is a non-zero element of
$\ZZ[X,B]$
\item $\Rightarrow$ there exist a positive real number $\gamma$ and a
non-zero polynomial $g \in \ZZ[X]$ such that for all finite fields $\FF$
and $x \in \FF^k$ with $g(x) \neq 0$ we have $|f(\{x\} \times \FF^n)|
\geq \gamma |\FF|^n$
\end{itemize}
\end{slide}

%%%%%%%%%%%%%%%%%%%%%%%%%%%%%%%%%%%%%%%%%%%%%%%%%%%%%%%%%%%%%%%%%%%%%%%%

\begin{slide}{Implications for our AG problem}
\begin{itemize}
\item let $X_{\mathrm{good}}:=\{x\in\F^k \,:\, g(x)\neq 0\}$
\item for all $x\in X_{\mathrm{good}}$ let $W_{\mathrm{good}}^x\subseteq\F^n$ to the set such that
\[
w\in W_{\mathrm{good}}^x \Rightarrow 1 \le \eta_x^w \le d
\]
where $d$ is some constant independent of the size of $\F$
\item we can implement quantum sampling
\[
V_x|S_w^x\>=|w\>
\]
for all $(x,w)$ with $x\in X_{\mathrm{good}}$ and $w\in W_{\mathrm{good}}^x$ because
\begin{itemize}
\item $S_w^x$ have constant size for all good pairs $(x,w)$ and
\item all solutions $b\in S_w^x$ can be enumerated efficiently 
\end{itemize}
\end{itemize}
\end{slide}

%%%%%%%%%%%%%%%%%%%%%%%%%%%%%%%%%%%%%%%%%%%%%%%%%%%%%%%%%%%%%%%%%%%%%%%%

\begin{slide}{Success of Approximate Measurement}
if choose the number of copies $k$ to be equal to $n$ (the maximal degree of the hidden polynomial) then
\begin{itemize}
\item $|X_{\mathrm{good}}|=\Omega(p^n)$
\item $|W_{\mathrm{good}}^x|=\Omega(p^n)$
\item the overall success probability is bounded from below as follows
\begin{eqnarray*}
P_{{\mathrm success}} & = &
\frac{1}{p^{3n}}
\sum_{x \in X_{\mathrm{good}}}
\left( \sum_{w \in W_{\mathrm{good}}^x} \sqrt{\eta^x_w} \right)^2 \\
& \approx & 
\frac{1}{p^{3n}} \,
\Omega(p^n) \, 
\Big(
\Omega(p^n) 
\Big)^2 = \Omega(1)
\end{eqnarray*}
\end{itemize}
\end{slide}

%%%%%%%%%%%%%%%%%%%%%%%%%%%%%%%%%%%%%%%%%%%%%%%%%%%%%%%%%%%%%%%%%%%%%%%%

\begin{slide}{Conclusions and Possible Extensions}
\begin{itemize}
\item definition of a new black-box problem (a special case of the hidden polynomial problem)
\item classically hard (polynomial classical query complexity) 
\item polylogarithmic quantum algorithm (for fixed degree)   
\item $\Rightarrow$ what useful real-life problem can be reduced to this black-box problem???
\item possible extensions:
\begin{itemize}
\item polynomials over non-prime fields; this is straightforward
\item polynomials over rings such as $\Z_d$???
\item multivariate polynomials??? some cases can be easily reduced to univariate case, e.g., hidden hyperbola
\end{itemize}
\end{itemize}
\end{slide}

%%%%%%%%%%%%%%%%%%%%%%%%%%%%%%%%%%%%%%%%%%%%%%%%%%%%%%%%%%%%%%%%%%%%%%%%

\end{document}